\documentclass[11pt]{article}
\usepackage{amsmath,amssymb,booktabs}
\usepackage[margin=1in]{geometry}
\begin{document}

\section*{Step 207: Candidate Transport-Sampling Comparison}

This step tests two candidate formulas for
\[
  \kappa_i(\tau)=\bigl[T_{1/2}^*K_{1/2}^{\Gamma}(\cdot,\rho_i)\bigr](1/2+i\tau)
\]
on the same 200-node Gauss--Legendre grid in \([-40,40]\).

\subsection*{CAND1: Burnol boundary kernel}

Burnol 2002, Theorem 8 states that \(B(E_\lambda)\) coincides isometrically
with the space of completed Mellin transforms of \(K_\lambda\).  With Burnol
2002 equation 1,
\[
 K_{1/2}^{\Gamma}(s,w)=
 \frac{E_{1/2}(s)\overline{E_{1/2}(w)}
       -E_{1/2}(1-s)\overline{E_{1/2}(1-w)}}{s+w-1},
\]
the first candidate is
\[
  \mathrm{CAND1}_i(\tau)=K_{1/2}^{\Gamma}(1/2+i\tau,\rho_i).
\]
The step uses the Step 205 Burnol-boundary grid computed from Burnol 2002
Theorem 8 and equation 1.

\subsection*{CAND2: zeta-dual candidate}

Burnol 2004 [19], Theorem 3.1 identifies the \(a=1\), simple-zero dual system
with inverse Mellin transforms of
\[
 \frac{\zeta(s)}{(s-\rho)\zeta'(\rho)\pi^{-\rho/2}\Gamma(\rho/2)}.
\]
For \(a=1/2\), Burnol states that suitable linear combinations are required.
The tested candidate is therefore the single-term normalization probe
\[
  \mathrm{CAND2}_i(\tau)=
  \frac{\zeta(1/2+i\tau)}
  {(1/2+i\tau-\rho_i)\zeta'(\rho_i)\pi^{-\rho_i/2}\Gamma(\rho_i/2)}.
\]

\subsection*{Ratio test}

The pointwise ratios \(\mathrm{CAND1}_i(\tau)/\mathrm{CAND2}_i(\tau)\)
would be constant in \(\tau\) if the candidates differed only by a
normalization.  The observed nominal ranges are:
\[
\begin{array}{c|c|c}
i & \min |\mathrm{CAND1}/\mathrm{CAND2}| & \max |\mathrm{CAND1}/\mathrm{CAND2}|\\
\hline
1 & 2.5008929629627379\cdot10^{-26} & 1.7748924779317304\cdot10^{-10}\\
2 & 7.4226939561589834\cdot10^{-30} & 2.6079509268902143\cdot10^{-15}\\
3 & 1.3962291433488764\cdot10^{-33} & 3.0769935351035648\cdot10^{-18}
\end{array}
\]
Thus the candidates do not agree up to a constant at the available numerical
level.  The imported Step 205 CAND1 error labels are conservative and dominate
the CAND1 values, so this is a transport-sampling obstruction rather than a
certified analytic inequality.

\subsection*{Consequence for \(\Xi_{\rm matrix\_source}\)}

The inherited compressed-HS formula is
\[
  \Xi_{\rm matrix\_source}=\operatorname{tr}(G^{-1}c^\dagger G^{-1}c),
  \qquad
  c_{ij}=\langle \kappa_i,(I-P_\infty)M_{m_\ell}P_\infty\kappa_j\rangle .
\]
Since the transport-sampling formula is unresolved, Step 207 does not select
either candidate and does not evaluate \(c_{ij}\) or \(\Xi_{\rm matrix\_source}\).

\subsection*{Verdict}

\[
  \boxed{\texttt{V\_kappa\_tau\_candidates\_disagree}.}
\]

The next typed dependency is the explicit action of
\[
  T_a^*=U_\infty J_a^*M_\Gamma^*
\]
on \(B(E_a)\) reproducing kernels, or a sharpened proof that the \(a<1\)
Burnol dual system reduces to a specified linear combination matching the
boundary kernel.

\end{document}
